\chapter*{Definitions}

For the purposes of this disclosure, the following terms shall have the meanings set forth below. These definitions are intended to provide clarity and shall not be construed to limit the scope of the invention unless explicitly stated.

\begin{definition}[Feedback-Coupled Agentic System (FCAS)]
A computational system in which telemetry outputs are not passive records but actively influence subsequent internal state transitions, including but not limited to model weights, policy updates, or decision functions.
\end{definition}

\begin{definition}[Telemetry]
Structured data generated during execution of a computational process, including observations, actions, rewards, or state transitions, which may be used for monitoring, training, or feedback.
\end{definition}

\begin{definition}[Deterministic State]
A system condition in which identical inputs, including telemetry disclosures, result in identical internal computational outcomes across executions and across heterogeneous hardware environments.
\end{definition}

\begin{definition}[Canonical Serialization]
A deterministic encoding process in which structured data is transformed into a byte sequence using a fixed and unambiguous ordering, representation, and precision, such that identical logical inputs always yield identical binary outputs.
\end{definition}

\begin{definition}[Fixed-Precision Scaling]
A numerical normalization process in which floating-point values are transformed into integer representations using a predefined scaling factor, thereby eliminating inconsistencies arising from floating-point arithmetic across architectures.
\end{definition}

\begin{definition}[Merkle-Chain (Lineage Chain)]
A cryptographic structure in which each data block is hashed and linked to the hash of a preceding block, forming a sequential chain that enables detection of any modification, omission, or reordering of historical data.
\end{definition}

\begin{definition}[Hash Function]
A deterministic algorithm that maps input data of arbitrary size to a fixed-size output, such that it is computationally infeasible to reverse or to find collisions under standard cryptographic assumptions.
\end{definition}

\begin{definition}[Proof of Reflexion (PoR)]
A cryptographic proof mechanism that binds a telemetry disclosure to the originating computational state or agent, ensuring that the data was generated by an authorized and consistent internal process.
\end{definition}

\begin{definition}[Threshold Signature Scheme]
A cryptographic protocol in which a valid signature is produced only when a predefined minimum number of independent participants (threshold $t$ out of $n$) contribute partial signatures.
\end{definition}

\begin{definition}[Zero-Trust Model]
A security model in which no component, node, or process is inherently trusted, and every interaction must be explicitly verified through cryptographic or formal validation mechanisms.
\end{definition}

\begin{definition}[Telemetry Batch]
A discrete unit of telemetry data grouped for processing, serialization, hashing, and verification within the system.
\end{definition}

\begin{definition}[Lineage Integrity]
The property that the complete sequence of telemetry batches is preserved without alteration, deletion, or reordering, and can be cryptographically verified end-to-end.
\end{definition}

\begin{definition}[Adversary]
A probabilistic computational entity capable of attempting to manipulate, forge, replay, or disrupt telemetry data, system state, or communication channels within defined computational bounds.
\end{definition}

\begin{definition}[Audit-Grade Verification]
A level of assurance in which system outputs and historical data can be independently verified for correctness, authenticity, and completeness using cryptographic and procedural methods.
\end{definition}

\begin{definition}[Ingestion]
The process by which telemetry data is received, validated, cryptographically anchored, and stored within a persistent system.
\end{definition}

\begin{definition}[Replay Attack]
An adversarial action in which previously valid data is retransmitted or reused to produce unintended or fraudulent system outcomes.
\end{definition}

\begin{definition}[Semantic Validity]
The property that telemetry data not only satisfies structural and cryptographic checks but also conforms to expected behavioral or domain-specific constraints of the system.
\end{definition}

\begin{definition}[Heterogeneous Architecture]
A computing environment composed of multiple hardware or software platforms with differing numerical, computational, or execution characteristics.
\end{definition}

\begin{definition}[Immutable Ledger]
A data structure in which recorded entries cannot be altered retroactively without detection through cryptographic verification.
\end{definition}

\begin{definition}[Sidecar Component]
An auxiliary process or service that operates alongside a primary application to provide additional functionality such as monitoring, logging, or cryptographic processing without modifying the core application logic.
\end{definition}